\begin{textsample}{POS Dim 5 – global\_north\_2023\_11 – Score 10.00 – global\_north\_2023\_11/103627217.txt}  \label{ex:f5_pos_415}
End To Killing Of Children In Armed Conflict, UN Committee Urges GENEVA ( 20 November 2023 ) -- With one out of every five children worldwide living within armed conflict zones, the UN Committee on the Rights of the Child marks World Children ' s Day in a sombre mood and calls for ceasefires and a return to basics of humanitarian \textbf{law} to \textbf{safeguard} all children.
The Committee today issued the following statement :"World Children ' s Day has generally been regarded as a day to celebrate the gains made since the United Nations adopted the \textbf{Convention} on the Rights of the Child in 1989.
Thirty-four years later today, it, however, has become a day for mourning for the many children who have recently died in armed conflict.
More than 4,600 children have been killed in Gaza in only five weeks.
This war has claimed the lives of more children in a shorter time and with a level of brutality that we have not witnessed in recent decades.
The Committee has previously urged for a ceasefire.
Unfortunately, the UN Security Council has not @ @ @ @ @ @ @ @ @ @ November 2023 \textbf{resolution} of the Security Council calling for humanitarian pauses and corridors is a positive step by the \textbf{international} community, it does not end the war that is waging on children -- it simply makes it possible for children to be saved from being killed on some days, but not on other days.
Advertisement - \textbf{scroll} to continue reading Are you getting our free newsletter?
There are 468 million children worldwide living in armed conflict zones, according to Save the Children ' s research, accounting for about 20 \% of the world ' s 2.4 billion children population, based on UNICEF ' s statistics.
On World Children ' s Day, the Committee also wants to underscore that while the armed conflict in the Occupied Palestinian Territory is at the forefront of our minds, we remain acutely concerned that thousands of children are dying in armed conflict in many parts of the world, including in Ukraine, Afghanistan, Yemen, Syria, Myanmar, Haiti, Sudan, Mali, Niger, Burkina Faso, Democratic Republic of the @ @ @ @ @ @ @ @ @ @, the global figure of children killed or maimed was 8,630.
Of deep concern is the fact that up to 4,000 children were denied humanitarian access last year.
Given the current situation in Gaza, the number of child victims of these grave human rights \textbf{violations} is rising exponentially.
The plight of girls affected by armed conflict is also at a crisis point.
In Sudan and Haiti, there are verified reports of abduction and rape of girls, and concerns have been raised by the Special Rapporteur on trafficking in persons, especially women and children about the deterioration in access to humanitarian services is driving girls towards being recruited by armed groups.
Children of so-called ' foreign fighters ' are a further area of concern.
The Committee has recommended in three complaints under its communications procedure that children in the camps in Northeast Syria should be repatriated.
While some States have \textbf{acted} to return children and their mothers, an estimated 31,000 children are still living in abysmal conditions in the camps.
The Committee also remains very concerned @ @ @ @ @ @ @ @ @ @ they reach early adolescence, as well as several hundred boys who are in prison.
The Committee \textbf{recognises} World Children ' s Day in a sombre mood.
In the face of wars affecting children around the globe, we call again for ceasefires, for a return to the basics of humanitarian \textbf{law}, and for thorough investigations by competent authorities of all grave \textbf{violations} against children in the context of armed conflict."Did you know Scoop has an Ethical Paywall?
If you ' re using Scoop for work, your organisation needs to pay a small license fee with Scoop Pro.
We think that ' s \textbf{fair}, because your organisation is benefiting from using our news resources.
In return, we ' ll also give your team access to pro news tools and keep Scoop free for personal use, because public access to news is important!
Among the key issues facing diplomats is securing the release of a reported 199 Israeli hostages, seized during the Hamas raid."History is watching,"says Emergency Relief Coordinator @ @ @ @ @ @ @ @ @ @ those hostages.
Of course, there ' s a history between Palestinian people and the Israeli people, and I ' m not denying any of that.
But that \textbf{act} alone lit a fire, which can only be put out with the release of those hostages."More Families in western Afghanistan are reeling after a fourth earthquake hit Herat Province, crumbling buildings and forcing people to flee once again, with thousands now living in tents exposed to fierce winds and dust storms.
The latest 6.3 magnitude earthquake hit 30 km outside of Herat on Sunday, shattering communities still reeling from strong and shallow aftershocks.
More UN Spokesperson St? phane Dujarric said some 1.1M people would be expected to leave northern Gaza and that such a movement would be"impossible"without devastating humanitarian consequences and appeals for the order to be rescinded.
The WHO joined the call for Israel to rescind the relocation order, which amounted to a"death sentence"for many.
More By October 10, reports indicated that fixed-line internet @ @ @ @ @ @ @ @ @ @ networks are all seriously compromised.
With significant and increasing damage to the electrical grid, orders by the Israeli Ministry of Energy to stop supplying electricity and the last remaining power station now out of fuel, many are no longer able to charge devices that are essential to communicate and access information.
More

% matched lemmas: act, convention, fair, international, law, recognise, resolution, safeguard, scroll, violation
\end{textsample}
